\documentclass[11pt, reqno]{amsart}

\usepackage{../ayud-template}
../LaTeX/general.tex

\usepackage{tikz}
\usetikzlibrary{babel,cd}

\usepackage{minted}
\usemintedstyle{colorful}
\setminted{
	mathescape,
	% linenos,
	autogobble,
	bgcolor=black!5,
	tabsize=2,
	fontsize=\small
}

% \usepackage{../ayud-template}
% ../LaTeX/general.tex
% % \input{../graphics.tex}

% \usepackage{multicol}
\DeclareMathOperator{\LT}{\text{\textsc{lt}}}

\title{Bases de Gröbner}
\date{\DTMdate{2025-11-05}}

\author{José Cuevas Barrientos}
\email{josecuevasbtos@uc.cl}
\urladdr{https://josecuevas.xyz/teach/2025-2-alg/}

\logo{../puc_negro.png}
\institution{Pontificia Universidad Católica de Chile y Universidad de Chile}
\department{Facultad de Matemáticas}
\course{Álgebra II}
\coursecode{MPG3201}
\professor{José Samper}

\begin{document}

\maketitle

\section{Teoría}
De aquí en adelante, $\N = \Z_{\ge 0}$.
Escribiremos una tupla $\vec\alpha = (\alpha_1, \dots, \alpha_n) \in \N^n$ en negritas y emplearemos notación multiíndice, es decir,
en $k[x_1, \dots, x_n] =: k[\vec x]$ se denota
\[
	\vec{x}^{\vec\alpha} := x_1^{\alpha_1} \cdots x_n^{\alpha_n}.
\]
\begin{mydef}
	Un \strong{orden monomial} $\preceq$ sobre $\N^n$ es una relación de orden total tal que:
	\begin{enumerate}[{OM}1., ref=OM\arabic*]
		\item Es un buen orden; es decir, todo subconjunto de $\N^n$ tiene $\preceq$-mínimo.
		\item\label{ax:sum_invariant}
			Si $\vec\alpha \prec \vec\beta$ y $\vec\gamma \in \N^n$, entonces $\vec\alpha + \vec\gamma \prec \vec\beta + \vec\gamma$.
	\end{enumerate}
\end{mydef}
\begin{obs}
	Sea $\preceq$ un orden total en $\N^n$ que satisface \ref{ax:sum_invariant}.
	Entonces $\preceq$ es un orden monomial syss $\vec 0 \in \N^n$ es el $\preceq$-mínimo de $\N^n$
	(cfr.\ \cite{cox:ideals}, Cor.~2.4.6).
\end{obs}
\begin{ex}
	Recuerde que el \emph{orden lexicográfico} $\leq_{\rm lex}$ (\textquote{del diccionario}) en $\N^n$ es aquel en donde $\vec\alpha
	<_{\rm lex} \vec\beta$ si, $\alpha_i < \beta_i$, donde $1 \le i \le n$ es la mínima coordenada $j$ tal que $\alpha_j \ne \beta_j$.
	Por ejemplo, $(10, 1, 2025) <_{\rm lex} (11, 0, 0)$.
	Este orden es monomial.

	Similarmente, definimos el orden $\leq_{\rm grlex}$ en $\N^n$ dado por
	\[
		(\alpha_1, \dots, \alpha_n) \leq_{\rm grlex} (\beta_1, \dots, \beta_n) \iff
		(\alpha_n, \dots, \alpha_1) \leq_{\rm lex} (\beta_n, \dots, \beta_1).
	\]
\end{ex}

\begin{thm}[de las bases de Hilbert]
	Todo ideal en $k[\vec x]$ está finitamente generado.
\end{thm}
Querremos ahora obtener un sistema generador para un ideal $\mathfrak{a}$ de manera algorítmica, replicando el caso de una variable.

\begin{mydef}
	Sea $f = \sum_{\vec\alpha} c_{\vec\alpha} \vec x^{\vec\alpha} \in k[\vec x]$ un polinomio.
	Tras fijar un orden monomial $\preceq$ sobre $\N^n$ se define su \emph{término lider} como $\LT(f) = c_{\vec\beta} \vec
	x^{\vec\beta}$, donde $\vec\beta$ es el $\preceq$-máximo multiíndice $\vec\alpha$ tal que $c_{\vec\alpha} \ne 0$.

	Sea $\mathfrak{a} \nsle k[x_1, \dots, x_n]$ un ideal.
	Se define
	\[
		\LT\mathfrak{a} := \sum_{f \in \mathfrak{a}} \LT(f)k[\vec x],
	\]
	es decir, es el ideal generado por los términos lider.
	Se dice que un subconjunto finito $\{ g_1, \dots, g_t \} \subseteq \mathfrak{a}$ es una \strong{base de Gröbner} si $\LT\mathfrak{a}
	= (\LT(g_1), \dots, \LT(g_t))$.
\end{mydef}
Siguiendo la misma demostración de que $k[x]$ (en una variable) es finitamente generado, obtenemos que:
\begin{cor}
	Toda base de Gröbner de un ideal $\mathfrak{a} \nsle k[\vec x]$ es un sistema generador.
\end{cor}

\begin{obs}
	Sea $\mathfrak{a} = (f_1, \dots, f_r)$, entonces no es cierto que $\LT\mathfrak{a} = (\LT(f_1), \dots, \LT(f_r))$, sino que solo hay
	una inclusión \textquote{$\supseteq$} en general.
	% Vea el siguiente 
\end{obs}
\begin{ex}\label{ex:non-grob}
	En $k[x, y]$ tome el orden ${\leq} := {\leq_{\rm lex}}$ (donde $y < x$),
	y considere el ideal $\mathfrak{a} = (f_1, f_2)$ con $f_1 = x^3 - 2xy$ y $f_2 = x^2y - 2y^2 + x$.
	Entonces $\LT(f_1) = x^3$ y $\LT(f_2) = x^2y$, de modo que $(\LT(f_1), \LT(f_2)) = x^2\cdot (x, y)$.

	No obstante,
	\[
		x f_1 - y f_2 = x^2 \in \mathfrak{a},
	\]
	y $\LT(x^2) = x^2 \notin (\LT(f_1), \LT(f_2))$.
	Así pues, $f_1, f_2$ \emph{no} es una ($\le_{\rm lex}$-)base de Gröbner de $\mathfrak{a}$.
	% Cambiando el orden por $\le_{\rm grlex}$ tampoco funciona pues
\end{ex}

\begin{mydef}
	Dada una tupla ordenada $F := (f_1, \dots, f_s)$ de polinomios en $k[\vec x]$ y otro $g \in k[\vec x]$, denotamos por
	$\overline{g}^F$ al resto de $g$ por división por $f_1, \dots, f_s$ en orden.
	Dados $f, g \in k[\vec x]$ (y tras fijar un orden monomial) denotamos por
	\[
		S(f, g) := \frac{\vec x^{\vec\gamma}}{\LT(f)} \cdot f - \frac{\vec x^{\vec\gamma}}{\LT(g)} \cdot g,
	\]
	donde $\gamma_j := \max\{ \alpha_j, \beta_j \}$ y donde $\vec x^{\vec\alpha} = \LT(f)$ y $\vec x^{\vec\beta} = \LT(g)$ salvo
	escalares.
\end{mydef}
\begin{thm}[algoritmo de Buchberger]
	Sea $G := (g_1, \dots, g_s)$ una tupla de polinomios en $k[\vec x]$ que genera el ideal $\mathfrak{a}$.
	% Tras adjuntar sucesivamente $\overline{S(p, q)}^G$ donde $p, q \in G$, podemos obtendremos una base de Gröbner del ideal generado
	% por $G$.
	Si $\overline{S(g_i, g_j)}^G = 0$ para todo $i \ne j$, entonces $G$ es una base de Gröbner.
\end{thm}
\begin{ex}
	Volvamos al ejemplo de $\mathfrak{a} = (f_1, f_2)$ con $f_1 = x^3 - 2xy$ y $f_2 = x^2y - 2y^2 + x$.
	Aquí,
	\[
		S(f_1, f_2) = yf_1 - xf_2 = -2xy^2 + 2xy^2 - x^2 = -x^2,
	\]
	y como $\LT(x^2) < \LT(f_1) < \LT(f_2)$, vemos que el resto $f_3$ es sí mismo.
	Luego, considere $G_1 := (f_1, f_2, f_3)$.

	Calculamos $S(f_1, f_3) = -2xy^2 =: f_4$ y $S(f_2, f_3) = -2y^2 + x =: f_5$, en ambos casos su resto es sí mismo.
	Finalmente el conjunto $G_2 := (f_1, \dots, f_5)$ es una base de Gröbner.
\end{ex}

\begin{additional}
\appendix
\printbibliography
\end{additional}

\end{document}
